% Generated draft: requires mathematical and source review.
\documentclass[11pt]{article}
\usepackage[margin=25mm]{geometry}
\usepackage{fontspec}
\setmainfont{Latin Modern Roman}
\usepackage{amsmath,amssymb,mathtools,unicode-math}
\setmathfont{Latin Modern Math}
\usepackage{enumitem,needspace,xurl,hyperref}
\setlength{\parindent}{0pt}
\setlength{\parskip}{5pt}
\setlength{\emergencystretch}{3em}
\newcommand{\N}{\mathbb{N}}
\newcommand{\Z}{\mathbb{Z}}
\newcommand{\Q}{\mathbb{Q}}
\newcommand{\R}{\mathbb{R}}
\newcommand{\C}{\mathbb{C}}
\newcommand{\F}{\mathbb{F}}
\newcommand{\A}{\mathbb{A}}
\newcommand{\PP}{\mathbb{P}}
\newcommand{\E}{\mathbb{E}}
\newcommand{\eps}{\varepsilon}
\newcommand{\dd}{\,\mathrm{d}}
\begin{document}
\section*{UnsolvedMath ID: 1200011}
\textbf{Source problem code:} catalogue entry
\subsection{Definitions and mathematical statement}
Under the usual connected undirected simple-graph interpretation, let $G$ be an infinite $d$-regular graph, $d\geq2$. Its simple random walk has transition operator
\[
(Pf)(v)=\frac1d\sum_{w\sim v}f(w)
\]
on $\ell^2(V(G))$. Call $G$ Ramanujan when its spectral radius satisfies $\|P\|\leq2\sqrt{d-1}/d$, the tree bound. Let $T_d$ be the infinite $d$-regular tree, rooted at $o$. For a rooted graph $(H,v)$, write $B_H(v,r)$ for the rooted induced subgraph on vertices of distance at most $r$ from $v$.

Fix any initial vertex $x\in V(G)$ and let $(X_n)_{n\geq0}$ be simple random walk with $X_0=x$. Ab\'ert's neighborhood-sampling question has the following local-distribution formulation:
\[
\forall r\in\Z_{\geq0},\qquad
\lim_{n\to\infty}\Pr_x\!\left[
(B_G(X_n,r),X_n)\cong(B_{T_d}(o,r),o)
\right]=1.
\]
The initial vertex is arbitrary, the radius is fixed before taking the limit, and rooted isomorphism is meant. No uniformity over growing radii is requested. The source's one-line formulation does not explicitly give graph conventions or a mode of convergence; this is the standard convergence-in-distribution interpretation of random-walk neighborhood sampling, not an assertion that every realized walk eventually sees only tree neighborhoods.

The supplied research summary reports an affirmative answer by Lyons--Peres (2015), Theorem 1.2: for every fixed $L\geq1$, the probability that $X_n$ lies on a nontrivial cycle of length at most $L$ tends to zero. It also reports almost-sure zero density of times traversing nontrivial cycles (Theorem 1.1) and exponential quantitative control (Theorem 4.2). These additional assertions involve the cited work's cycle/traversal conventions, which are not reproduced here. The summary emphasizes that no unimodularity, stationarity, or transitivity hypothesis is needed, and describes the proof as comparing simple and nonbacktracking walks to control spectral radius through cogrowth.
\subsection{Short English statement}
At each fixed radius, does the neighborhood around a late position of simple random walk on an infinite regular Ramanujan graph look like the regular tree with probability tending to one? The supplied summary reports yes, without assuming stationarity, unimodularity, or transitivity.
\subsection{Sources}
\begin{itemize}
\item Local UnsolvedMath registry, permanent ID 1200011, exact source code catalogue entry, ``Some Questions --- Question 11''; status solved, classification SOLVED-IN-LITERATURE.
\item Ab\'ert, \emph{Some questions} (2010), Question 11, \url{https://www.renyi.hu/~abert/questions.pdf}. The provenance records manual transcription from a public author-hosted PDF, accessed July 29, 2026, and says the 2010 list was inspected but current status was not independently verified.
\item Lyons--Peres (2015), Theorems 1.2, 1.1, and 4.2, as cited in the supplied research summary for short-cycle probability decay, almost-sure traversal-density zero, and exponential quantitative control, respectively. A title, publication details, and exact cycle definitions are not supplied.
\end{itemize}
\end{document}
